% Sablon pentru realizarea lucrarii de licenta, conform cu recomandarile
% din ghidul de redactare:
% - https://fmi.unibuc.ro/finalizare-studii/

% Multumiri lui Gabriel Majeri, acest sablon a fost creat pe baza
% codului sursa a lucrarii sale de licenta. 
% Codul sursa: https://github.com/GabrielMajeri/bachelors-thesis
% Website: https://www.gabrielmajeri.ro/
%
% Acest sablon este licentiat sub Creative Commons Attribution 4.0 International License.
% modificat 8 noiembrie 2024

\documentclass[12pt, a4paper]{article}

% Suport pentru diacritice și alte simboluri
\usepackage{fontspec}

% Font Times New Roman
\usepackage{times}

% Suport pentru mai multe limbi
\usepackage{polyglossia}

% Setează limba textului la română
\setdefaultlanguage{romanian}
% Am nevoie de engleză pentru rezumat
\setotherlanguages{english}

% Indentează și primul paragraf al fiecărei noi secțiuni
\SetLanguageKeys{romanian}{indentfirst=true}

% Suport pentru diferite stiluri de ghilimele
\usepackage{csquotes}

\DeclareQuoteStyle{romanian}
  {\quotedblbase}
  {\textquotedblright}
  {\guillemotleft}
  {\guillemotright}

% Utilizează biblatex pentru referințe bibliografice
\usepackage[
    maxbibnames=50,
    sorting=nty
]{biblatex}

\addbibresource{bibliography.bib}

% Setează spațiere inter-linie (doublespacing din LaTeX este echivalentul setarii 1.5 in Microsoft Word)
\usepackage{setspace}
\doublespacing

% Modificarea geometriei paginii (cu margini de 2,5 cm) 
\usepackage[margin=2.5cm]{geometry}

% Include funcțiile de grafică
\usepackage{graphicx}
% Încarcă imaginile din directorul `images`
\graphicspath{{./images/}}

% Listări de cod
\usepackage{listings}

% Linkuri interactive în PDF
\usepackage[
    colorlinks,
    linkcolor={black},
    menucolor={black},
    citecolor={black},
    urlcolor={blue}
]{hyperref}

% Simboluri matematice codificate Unicode
\usepackage[warnings-off={mathtools-colon,mathtools-overbracket}]{unicode-math}

% Comenzi matematice
\usepackage{amsmath}
\usepackage{mathtools}

% Formule matematice
\newcommand{\bigO}[1]{\symcal{O}\left(#1\right)}
\DeclarePairedDelimiter\abs{\lvert}{\rvert}

% Suport pentru rezumat în două limbi
% Bazat pe https://tex.stackexchange.com/a/70818
\newenvironment{abstractpage}
  {\cleardoublepage\vspace*{\fill}\thispagestyle{empty}}
  {\vfill\cleardoublepage}
\renewenvironment{abstract}[1]
  {\bigskip\selectlanguage{#1}%
   \begin{center}\bfseries\abstractname\end{center}}
  {\par\bigskip}

% Suport pentru anexe
\usepackage{appendix}

% Stiluri diferite de headere și footere
\usepackage{fancyhdr}

% Metadate
\title{Titlul lucrării de licență}
\author{Numele studentului}

% Generează variabilele cu @
\makeatletter

\begin{document}

% Front matter
\cleardoublepage
\let\ps@plain

% Pagina de titlu
\begin{titlepage}

% Redu marginile
\newgeometry{left=2cm,right=2cm,bottom=1cm}

\begin{figure}[!htb]
    \centering
    \begin{minipage}{0.2\textwidth}
        \includegraphics[width=\linewidth]{logo-ub.png}
    \end{minipage}
    \begin{minipage}{0.5\textwidth}
        \large
        \vspace{0.2cm}
        \begin{center}
            \textbf{UNIVERSITATEA DIN BUCUREȘTI}
        \end{center}
        \vspace{0.3cm}
        \begin{center}
            \textbf{
                FACULTATEA DE \\
                MATEMATICĂ ȘI INFORMATICĂ
            }
        \end{center}
    \end{minipage}
    \begin{minipage}{0.2\textwidth}
        \includegraphics[width=\linewidth]{logo-fmi.png}
    \end{minipage}
\end{figure}

\begin{center}
\textbf{SPECIALIZAREA INFORMATICĂ}
\end{center}

\vspace{1cm}

\begin{center}
\Large \textbf{Lucrare de licență}
\end{center}

\begin{center}
\huge \textbf{\MakeUppercase{\@title}}
\end{center}

\vspace{3cm}

\begin{center}
\large \textbf{Absolvent \\ \@author}
\end{center}

\vspace{0.25cm}

\begin{center}
\large \textbf{Coordonator științific \\ Titlul și numele profesorului coordonatorului}
\end{center}

\vspace{2cm}

\begin{center}
\Large \textbf{București, iunie 2021}
\end{center}
\end{titlepage}
\restoregeometry
\newgeometry{
    margin=2.5cm
}

\fancypagestyle{main}{
  \fancyhf{}
  \renewcommand\headrulewidth{0pt}
  \fancyhead[C]{}
  \fancyfoot[C]{\thepage}
}

\addtocounter{page}{1}

% Rezumatul
\begin{abstractpage}

\begin{abstract}{romanian}
Lorem ipsum dolor sit amet, consectetur adipiscing elit. Fusce vitae eros sit amet sem ornare varius. Duis eget felis eget risus posuere luctus. Integer odio metus, eleifend at nunc vitae, rutrum fermentum leo. Quisque rutrum vitae risus nec porta. Nunc eu orci euismod, ornare risus at, accumsan augue. Ut tincidunt pharetra convallis. Maecenas ut pretium ex. Morbi tellus dui, viverra quis augue at, tincidunt hendrerit orci. Lorem ipsum dolor sit amet, consectetur adipiscing elit. Aliquam quis sollicitudin nunc. Sed sollicitudin purus dapibus mi fringilla, nec tincidunt nunc eleifend. Nam ut molestie erat. Integer eros dolor, viverra quis massa at, auctor.
\end{abstract}

\begin{abstract}{english}
Lorem ipsum dolor sit amet, consectetur adipiscing elit. Fusce vitae eros sit amet sem ornare varius. Duis eget felis eget risus posuere luctus. Integer odio metus, eleifend at nunc vitae, rutrum fermentum leo. Quisque rutrum vitae risus nec porta. Nunc eu orci euismod, ornare risus at, accumsan augue. Ut tincidunt pharetra convallis. Maecenas ut pretium ex. Morbi tellus dui, viverra quis augue at, tincidunt hendrerit orci. Lorem ipsum dolor sit amet, consectetur adipiscing elit. Aliquam quis sollicitudin nunc. Sed sollicitudin purus dapibus mi fringilla, nec tincidunt nunc eleifend. Nam ut molestie erat. Integer eros dolor, viverra quis massa at, auctor.
\end{abstract}

\end{abstractpage}

\tableofcontents

% Main matter
\cleardoublepage
\pagestyle{main}
\let\ps@plain\ps@main

\chapter{Introducere}

\section{Lorem ipsum}

Lorem ipsum dolor sit amet \cite{fake_citation}, consectetur adipiscing elit. Nam convallis dapibus ante in interdum. Aliquam vitae elit eu elit gravida facilisis. Proin tincidunt sodales elit vitae ornare. Suspendisse quis rutrum ipsum, eget tempor felis. Nulla facilisi. Vivamus mollis quam a est aliquam congue. Nulla risus turpis, rutrum ac lacus molestie, vestibulum ultricies enim. Suspendisse a urna ultrices quam pulvinar finibus ac vitae felis. Ut sit amet interdum lacus. Sed in cursus arcu. Vivamus quis eros efficitur, ullamcorper dui ut, blandit nisl. Interdum et malesuada fames ac ante ipsum primis in faucibus. Nunc interdum venenatis mi, sed placerat elit feugiat in. Praesent dignissim ante nec rutrum tempus.

Aenean accumsan nunc nunc, nec maximus ex ultricies et. In porta sollicitudin augue. Curabitur consectetur elit nec eros laoreet fermentum. Vivamus nec nisl a lectus scelerisque viverra at id tortor. Fusce et erat dictum, dictum odio nec, aliquam orci. Vestibulum gravida feugiat purus. Mauris non neque est. Sed commodo vel enim vel tincidunt. Sed faucibus dolor id dui dictum pellentesque. Vestibulum ante ipsum primis in faucibus orci luctus et ultrices posuere cubilia curae; Mauris interdum tempor purus ut tempor. Maecenas ac suscipit ex. Sed vitae diam nec neque ultricies faucibus. In at purus elit. Nunc placerat nunc tortor.

Fusce fermentum quis tellus vitae congue. Interdum et malesuada fames ac ante ipsum primis in faucibus. Proin convallis congue sapien semper volutpat. Curabitur rhoncus et nunc quis tempor. Quisque sollicitudin dolor id varius luctus. Donec sapien dolor, eleifend quis justo id, euismod ultrices nulla. Suspendisse leo quam, luctus in velit ac, venenatis molestie odio. Aliquam vitae feugiat elit. Sed auctor gravida leo, ut dictum enim tempus non. Suspendisse in dapibus urna, vitae pretium ligula. Morbi ut diam laoreet, convallis leo non, semper arcu. Sed vel lorem tortor. In auctor a diam sit amet convallis. Etiam ipsum ligula, luctus quis interdum ac, consequat vitae orci. Maecenas hendrerit turpis tristique, tincidunt nisi ut, mattis augue.

\section{Dolor sit amet}

Fusce fermentum quis tellus vitae congue. Interdum et malesuada fames ac ante ipsum primis in faucibus. Proin convallis congue sapien semper volutpat. Curabitur rhoncus et nunc quis tempor. Quisque sollicitudin dolor id varius luctus. Donec sapien dolor, eleifend quis justo id, euismod ultrices nulla. Suspendisse leo quam, luctus in velit ac, venenatis molestie odio. Aliquam vitae feugiat elit. Sed auctor gravida leo, ut dictum enim tempus non. Suspendisse in dapibus urna, vitae pretium ligula. Morbi ut diam laoreet, convallis leo non, semper arcu. Sed vel lorem tortor. In auctor a diam sit amet convallis. Etiam ipsum ligula, luctus quis interdum ac, consequat vitae orci. Maecenas hendrerit turpis tristique, tincidunt nisi ut, mattis augue.

\subsection{Quis tellus vitae}

Aenean accumsan nunc nunc, nec maximus ex ultricies et. In porta sollicitudin augue. Curabitur consectetur elit nec eros laoreet fermentum. Vivamus nec nisl a lectus scelerisque viverra at id tortor. Fusce et erat dictum, dictum odio nec, aliquam orci. Vestibulum gravida feugiat purus. Mauris non neque est. Sed commodo vel enim vel tincidunt. Sed faucibus dolor id dui dictum pellentesque. Vestibulum ante ipsum primis in faucibus orci luctus et ultrices posuere cubilia curae; Mauris interdum tempor purus ut tempor. Maecenas ac suscipit ex. Sed vitae diam nec neque ultricies faucibus. In at purus elit. Nunc placerat nunc tortor.
\chapter{Preliminarii}

Lorem ipsum dolor sit amet, consectetur adipiscing elit. Aenean dignissim metus justo, nec pharetra mauris tincidunt id. Praesent semper turpis quis faucibus pulvinar. Fusce ut justo nisi. Praesent vehicula blandit erat, sed dignissim justo bibendum lobortis. Vivamus fringilla, elit at pulvinar imperdiet, dui elit lobortis sapien, et vehicula urna ex et velit. Sed efficitur, neque sed egestas lobortis, diam leo pellentesque sem, nec gravida est nunc efficitur orci. Maecenas bibendum pharetra sapien, quis porttitor justo vestibulum a. Nunc tempus erat sed augue blandit, eu scelerisque dolor lobortis. Vestibulum quam arcu, malesuada quis felis eleifend, iaculis gravida massa. Vestibulum lacinia arcu nec risus laoreet porttitor. Quisque ut nisl consequat, sodales nisl id, pretium neque. Interdum et malesuada fames ac ante ipsum primis in faucibus. Vivamus tincidunt, ligula in fringilla laoreet, urna risus efficitur orci, eget pellentesque metus magna interdum lorem.

Quisque et ligula erat. Aliquam eget fringilla tortor. Nulla maximus, massa ac consectetur fringilla, ante velit porttitor justo, et euismod magna eros consequat nunc. Praesent lacinia nulla dui, lobortis mattis odio ullamcorper et. Proin vehicula massa in efficitur eleifend. Mauris ornare mi ac rutrum porta. Morbi sed magna vel sem egestas malesuada ac at turpis. Aenean eleifend pharetra massa, eget porta ligula eleifend ultricies. Donec vel imperdiet leo. Nullam commodo metus metus, sed semper velit ornare posuere.

Aliquam erat volutpat. Pellentesque vulputate massa sed semper lacinia. Duis hendrerit dolor et blandit malesuada. Curabitur posuere tellus at lacinia scelerisque. Donec hendrerit semper ullamcorper. Morbi ut semper metus. Aliquam pharetra sagittis dolor, non pulvinar purus euismod non. Phasellus condimentum a est tristique convallis. Nam ac mauris bibendum, lobortis ex at, rhoncus nibh. Fusce neque nisi, dignissim porttitor lorem sed, lobortis vehicula felis.
\chapter{Concluzii}

\printbibliography[heading=bibintoc]

\end{document}